% ============================================================
% CAPÍTULO 4: PROTEÇÃO JURÍDICA DA PERSONALIDADE 
%              NO DIREITO BRASILEIRO
% ============================================================

\chapter{Proteção Jurídica da Personalidade no Direito Brasileiro}
\label{cap:protecao-juridica}

\section{Introdução ao Capítulo}
\label{sec:intro-cap4}

O ordenamento jurídico brasileiro dispõe de um conjunto significativo de instrumentos normativos para a proteção da personalidade humana. A Constituição Federal de 1988, o Código Civil, a Lei Geral de Proteção de Dados Pessoais, o Marco Civil da Internet e a jurisprudência dos tribunais superiores compõem um arcabouço que, embora não tenha sido concebido especificamente para enfrentar os desafios da inteligência artificial, oferece categorias e princípios aplicáveis a esse novo contexto.

Este capítulo examina esses instrumentos, avaliando suas potencialidades e insuficiências para a proteção da personalidade na era da IA. O percurso é o seguinte: parte-se dos direitos da personalidade no Código Civil (\ref{sec:direitos-personalidade-cc}), analisa-se o tratamento constitucional da dignidade humana (\ref{sec:dignidade-constitucional}), examina-se a LGPD como marco regulatório da autodeterminação informativa (\ref{sec:lgpd}), avalia-se o Marco Civil da Internet (\ref{sec:mci}), e finalmente examina-se o Projeto de Lei nº 2.338/2023 como proposta de regulação específica da IA no Brasil (\ref{sec:pl-2338}).

\section{Os Direitos da Personalidade no Código Civil}
\label{sec:direitos-personalidade-cc}

\subsection{Fundamentos e Características}
\label{subsec:fundamentos-direitos-personalidade}

O Código Civil de 2002 dedicou, pela primeira vez na história legislativa brasileira, um capítulo autônomo aos direitos da personalidade (arts. 11 a 21). Essa inovação representou um avanço significativo na proteção jurídica da pessoa humana, superando a concepção meramente patrimonialista que predominava no Código Civil de 1916.

Os direitos da personalidade são definidos pela doutrina como aqueles direitos subjetivos que têm por objeto os atributos físicos, psíquicos e morais da pessoa, tutelando sua integridade física, intelectual e moral \cite{Amaral2003DBC}. Suas características fundamentais são a intransmissibilidade, a irrenunciabilidade, a indisponibilidade e a imprescritibilidade. O art. 11 do CC estabelece que ``com exceção dos casos previstos em lei, os direitos da personalidade são intransmissíveis e irrenunciáveis, não podendo o seu exercício sofrer limitação voluntária'' \cite{Brasil2002CC}.

O rol de direitos da personalidade no CC é exemplificativo, não taxativo. Estão expressamente previstos: a proteção ao nome (art. 16 a 19), à imagem (art. 20), à privacidade (art. 21), à integridade física e psíquica, e à honra. A doutrina reconhece, contudo, a existência de direitos da personalidade não expressamente codificados, como o direito à identidade pessoal, à autodeterminação informativa e ao esquecimento (este último objeto de intenso debate jurisprudencial).

\subsection{A Insuficiência do Modelo Civilista para os Desafios da IA}
\label{subsec:insuficiencia-civilista}

Embora os direitos da personalidade previstos no CC constituam um ponto de partida importante, sua concepção clássica mostra-se insuficiente para enfrentar os desafios específicos postos pela IA. Essa insuficiência manifesta-se em pelo menos três planos.

Em \textbf{primeiro lugar}, o modelo civilista tradicional é essencialmente reativo e reparatório: pressupõe a ocorrência de um dano concreto para que o titular possa exercer seu direito de exigir reparação. A IA, contudo, produz frequentemente danos difusos e sistêmicos que não se enquadram facilmente no esquema tradicional da responsabilidade civil (dano, nexo causal, culpa). A discriminação algorítmica, por exemplo, afeta grupos inteiros de forma indireta, tornando difícil a individualização do dano e a identificação do nexo causal.

Em \textbf{segundo lugar}, o direito da personalidade clássico foi concebido para tutelar a pessoa contra agressões identificáveis e localizáveis. A opacidade dos sistemas algorítmicos torna difícil sequer saber se e quando uma violação ocorreu. Como exercer o direito de exigir reparação se a própria existência do dano é encoberta pela caixa-preta algorítmica?

Em \textbf{terceiro lugar}, os direitos da personalidade foram pensados primordialmente como direitos de defesa contra o Estado (dimensão vertical), mas os principais riscos decorrentes da IA emanam de atores privados, grandes corporações de tecnologia, que operam em escala global e com poder econômico frequentemente superior ao de muitos Estados nacionais. Essa dimensão horizontal dos direitos fundamentais (eficácia privada ou horizontal) exige uma reconstrução teórica dos instrumentos de proteção.

\section{A Dignidade da Pessoa Humana como Fundamento Constitucional}
\label{sec:dignidade-constitucional}

\subsection{A Dignidade como Princípio Estruturante}
\label{subsec:dignidade-principio}

A Constituição Federal de 1988 elegeu a dignidade da pessoa humana como fundamento da República (art. 1º, III) e como princípio informador de todo o ordenamento jurídico. A doutrina constitucional brasileira, especialmente a partir das contribuições de Ingo Sarlet e Daniel Sarmento, tem desenvolvido um conceito de dignidade que vai além de sua formulação meramente negativa (proibição de tratamento degradante) para abranger também uma dimensão positiva de promoção das condições para o pleno desenvolvimento da personalidade \cite{Sarlet2021DF, Sarmento2016DPF}.

A dignidade da pessoa humana, nessa perspectiva, não é um direito entre outros, mas o fundamento axiológico que dá unidade e sentido ao sistema de direitos fundamentais. Ela opera como limite (proibindo intervenções degradantes), como dever de proteção (obrigando o Estado a proteger a pessoa contra agressões de terceiros) e como mandato de promoção (exigindo políticas públicas que criem as condições materiais para o desenvolvimento da personalidade).

Essas três dimensões são relevantes para a regulação da IA. Como \textbf{limite}, a dignidade proíbe a utilização de pessoas como meros meios para fins econômicos ou tecnológicos, vedando práticas como a manipulação algorítmica não consentida ou a redução da pessoa a um perfil de dados. Como \textbf{dever de proteção}, ela obriga o Estado a editar normas que protejam as pessoas contra os riscos da IA e a fiscalizar sua observância. Como \textbf{mandato de promoção}, ela exige políticas de inclusão digital, letramento algorítmico e acesso equitativo às tecnologias.

\subsection{A Dignidade e a Proteção Contra Riscos Algorítmicos}
\label{subsec:dignidade-riscos}

A aplicação do princípio da dignidade à proteção contra riscos algorítmicos encontra amparo na jurisprudência do Supremo Tribunal Federal. No julgamento da ADI nº 6.640 (que tratou da validade constitucional da LGPD), o STF reconheceu que a proteção de dados pessoais é expressão do direito fundamental à privacidade e à autodeterminação informativa, ambos ancorados no princípio da dignidade da pessoa humana \cite{STF2020ADI6640}.

A \MH{} oferece uma contribuição relevante a esse debate ao articular o princípio da dignidade com a crítica ao paradigma tecnocrático. Para a Encíclica, a dignidade não é apenas um conceito jurídico abstrato, mas uma realidade concreta que pode ser violada sempre que a pessoa for tratada como objeto, como meio ou como dado. A proteção contra a redução algorítmica da pessoa a um perfil ou a uma classificação é, portanto, uma exigência da própria dignidade \cite{LeaoXIV2026MH}.

\section{A Lei Geral de Proteção de Dados Pessoais}
\label{sec:lgpd}

\subsection{Estrutura e Princípios da LGPD}
\label{subsec:estrutura-lgpd}

A Lei nº 13.709/2018, Lei Geral de Proteção de Dados Pessoais, representa o marco normativo mais relevante para a proteção da personalidade na era digital no Brasil. Inspirada no Regulamento Geral de Proteção de Dados da União Europeia (GDPR), a LGPD estabelece um conjunto abrangente de regras para o tratamento de dados pessoais por pessoas naturais e jurídicas, públicas e privadas.

A LGPD se estrutura em torno de princípios fundamentais que orientam toda a atividade de tratamento de dados: finalidade, adequação, necessidade, livre acesso, qualidade dos dados, transparência, segurança, prevenção, não discriminação e responsabilização (art. 6º). Esses princípios funcionam como cláusulas gerais que devem informar a interpretação e aplicação de todas as disposições da Lei.

Para os fins deste trabalho, destacam-se três princípios especialmente relevantes:

O \textbf{princípio da necessidade} (art. 6º, III) determina que o tratamento de dados deve limitar-se ao mínimo necessário para a realização de suas finalidades. Esse princípio é particularmente importante para conter a coleta massiva de dados que alimenta os sistemas de IA, estabelecendo um limite à voracidade informacional do capitalismo de vigilância.

O \textbf{princípio da não discriminação} (art. 6º, IX) proíbe a utilização de dados pessoais para fins discriminatórios ilícitos ou abusivos. Esse princípio constitui um fundamento normativo direto para o combate à discriminação algorítmica.

O \textbf{princípio da responsabilização e prestação de contas} (art. 6º, X) exige que o controlador demonstre a adoção de medidas eficazes capazes de comprovar a observância das normas de proteção de dados. Esse princípio fundamenta a exigência de auditorias e avaliações de impacto à proteção de dados, essenciais para o controle preventivo de riscos algorítmicos.

\subsection{O Direito à Revisão de Decisões Automatizadas}
\label{subsec:revisao-automatizada}

O art. 20 da LGPD consagra o direito do titular de dados pessoais ``a solicitar revisão de decisões tomadas exclusivamente com base em tratamento automatizado de dados pessoais que afetem seus interesses, incluídas as decisões destinadas a definir o seu perfil pessoal, profissional, de consumo e de crédito ou os aspectos de sua personalidade''.

Esse dispositivo é inovador no direito brasileiro, pois reconhece que decisões algorítmicas podem afetar a personalidade humana e cria um mecanismo processual para sua contestação. O §1º do mesmo artigo determina que a revisão deve ser realizada por pessoa natural, estabelecendo a garantia de supervisão humana como contrapeso à automação decisória.

Contudo, a implementação prática desse direito enfrenta desafios significativos. Estudos apontam que o exercício do direito de revisão é frequentemente dificultado pela falta de transparência dos sistemas algorítmicos, pela assimetria de informações entre controlador e titular, e pela ausência de padrões claros para a realização da revisão \cite{Bioni2021PPD, Doneda2022DPD}. Além disso, a redação do art. 20 refere-se a decisões ``tomadas exclusivamente com base em tratamento automatizado'', o que exige um elevado standard probatório por parte do titular que deseja contestar a decisão.

\subsection{As Avaliações de Impacto Algorítmico}
\label{subsec:avaliacoes-impacto}

Um dos instrumentos mais promissores para a proteção preventiva contra riscos algorítmicos é a Avaliação de Impacto à Proteção de Dados Pessoais (AIPD), prevista no art. 38 da LGPD. A AIPD é um processo sistemático de identificação, análise e mitigação de riscos à privacidade e à proteção de dados associados a operações de tratamento.

A Autoridade Nacional de Proteção de Dados (ANPD) tem incentivado a realização de AIPDs especialmente para sistemas de IA de alto risco. A Resolução ANPD nº 4/2024 estabeleceu diretrizes para a elaboração de relatórios de impacto, incluindo a descrição do funcionamento do sistema, a identificação de riscos e as medidas de mitigação adotadas.

A doutrina tem proposto a ampliação desse instrumento para além da mera proteção de dados, transformando-o em uma verdadeira Avaliação de Impacto Algorítmico (AIA) que considere não apenas riscos à privacidade, mas também riscos de discriminação, exclusão, manipulação e violação de direitos fundamentais \cite{Doneda2022DPD, Edwards2017PC}. A \MH{} oferece fundamentos éticos para essa ampliação, ao insistir que a proteção da pessoa humana deve ser abrangente e não restrita a uma única dimensão de sua personalidade \cite{LeaoXIV2026MH}.

\section{O Marco Civil da Internet}
\label{sec:mci}

\subsection{Princípios e Direitos}
\label{subsec:principios-mci}

A Lei nº 12.965/2014, Marco Civil da Internet (MCI), estabelece princípios, garantias, direitos e deveres para o uso da Internet no Brasil. Embora não trate especificamente de IA, o MCI oferece fundamentos relevantes para a proteção da personalidade no ambiente digital.

O MCI consagra, em seu art. 3º, os princípios de neutralidade da rede, privacidade, proteção de dados pessoais, liberdade de expressão e responsabilização dos agentes. Esses princípios informam a interpretação de todo o ordenamento digital e devem orientar também a regulação da IA.

Especialmente relevante é o art. 10 do MCI, que determina a proteção dos registros de conexão e de acesso a aplicações de internet, ressalvadas as hipóteses de guarda obrigatória e de requisição judicial. Esse dispositivo estabelece um regime de proteção de dados que se aplica também aos dados utilizados por sistemas de IA.

\subsection{A Responsabilidade Civil dos Provedores}
\label{subsec:responsabilidade-mci}

O MCI estabelece um regime especial de responsabilidade civil para provedores de aplicação. O art. 19 determina que o provedor só pode ser responsabilizado por danos decorrentes de conteúdo gerado por terceiros se, após ordem judicial específica, não tomar as providências para tornar indisponível o conteúdo. Esse regime tem sido objeto de intenso debate no Supremo Tribunal Federal (RE nº 1.037.396/SP e RE nº 1.057.258/SP), que discute sua constitucionalidade e extensão.

Para os sistemas de IA gerativa, o regime do art. 19 mostra-se particularmente problemático. Se um modelo de IA produz conteúdo danoso (deepfake difamatório, informação falsa, discurso de ódio), aplica-se a responsabilidade do provedor? A resposta não é clara, pois o conteúdo não é propriamente ``gerado por terceiros'' (o usuário que solicitou a geração) nem pelo provedor (que apenas disponibiliza o modelo). Essa lacuna normativa evidencia a necessidade de um marco regulatório específico para a IA.

\section{O Projeto de Lei nº 2.338/2023}
\label{sec:pl-2338}

\subsection{Contexto e Conteúdo da Proposta}
\label{subsec:contexto-pl}

O Projeto de Lei nº 2.338/2023, em tramitação no Congresso Nacional, propõe um marco regulatório abrangente para a inteligência artificial no Brasil. Inspirado no \textit{AI Act} da União Europeia, o PL adota uma abordagem baseada em risco, classificando os sistemas de IA em diferentes níveis de risco (risco excessivo, alto risco, risco limitado e risco mínimo) e estabelecendo obrigações proporcionais ao nível de risco.

Para os fins deste trabalho, destacam-se três aspectos do PL.

Em primeiro lugar, o PL proíbe determinadas práticas consideradas de risco excessivo, como sistemas de pontuação social, reconhecimento emocional em ambientes de trabalho e educacionais, e categorização biométrica baseada em características sensíveis. Essas proibições alinham-se com as preocupações da \MH{} sobre a instrumentalização da pessoa humana pela tecnologia.

Em segundo lugar, o PL estabelece um conjunto de obrigações para sistemas de IA de alto risco, incluindo avaliação de impacto algorítmico, transparência, supervisão humana, documentação técnica e governança de dados. Essas obrigações concretizam, no plano normativo, muitos dos princípios éticos propostos pela \MH{}, especialmente transparência e responsabilidade.

Em terceiro lugar, o PL cria o Sistema Nacional de Regulação de Inteligência Artificial (SINRIA) como estrutura de governança com competência para expedir normas, fiscalizar e sancionar infrações. A eficácia do PL dependerá crucialmente da capacidade do SINRIA de exercer essas funções de forma independente, técnica e eficiente.

\subsection{Diálogo com a \MH{}}
\label{subsec:dialogo-pl-mh}

O PL 2.338/2023 e a \MH{} operam em registros distintos (o primeiro no âmbito da regulação técnica e jurídica, a segunda no campo da reflexão ética e antropológica), mas ambos convergem na preocupação fundamental com a proteção da pessoa humana diante dos riscos da IA.

A \MH{} oferece ao debate regulatório brasileiro uma fundamentação antropológica que o PL, por sua própria natureza técnica, não desenvolve. Enquanto o PL estabelece o que deve ser regulado (classificação de riscos) e como (obrigações, sanções, governança), a Encíclica explica por que a proteção da pessoa é necessária e quais valores estão em jogo.

Essa complementaridade é particularmente evidente em três pontos. Primeiro, a classificação de riscos adotada pelo PL encontra correspondência na gradação ética proposta pela \MH{}, que distingue entre usos aceitáveis, problemáticos e inaceitáveis da IA. Segundo, as obrigações de transparência e supervisão humana do PL ecoam a exigência ética de que a pessoa não seja submetida a decisões que não possa compreender ou controlar. Terceiro, a estrutura de governança do PL pode ser vista como uma concretização institucional do princípio de subsidiariedade, na medida em que estabelece uma instância pública de coordenação e fiscalização sem eliminar a responsabilidade dos atores privados.

\section{Síntese do Capítulo}
\label{sec:sumario-cap4}

O exame dos instrumentos jurídicos brasileiros para a proteção da personalidade revela um panorama misto. Por um lado, o ordenamento dispõe de fundamentos constitucionais sólidos (dignidade da pessoa humana), institutos civilistas consolidados (direitos da personalidade), um marco regulatório de proteção de dados em desenvolvimento (LGPD) e uma proposta de regulação específica da IA (PL 2.338/2023). Por outro lado, identificam-se insuficiências significativas:

\begin{enumerate}
    \item O \textbf{modelo civilista clássico}, de caráter reativo e reparatório, mostra-se inadequado para lidar com danos difusos e sistêmicos produzidos por sistemas algorítmicos.
    \item A \textbf{LGPD}, apesar de seu avanço, carece de instrumentos específicos para o controle da discriminação algorítmica e para a garantia efetiva do direito à explicação.
    \item O \textbf{MCI} não oferece resposta adequada para a responsabilidade civil por conteúdos gerados por IA.
    \item O \textbf{PL 2.338/2023}, embora promissor, ainda está em tramitação e seu texto final é incerto.
\end{enumerate}

Diante desse quadro, o capítulo seguinte propõe-se a extrair das contribuições da \MH{} orientações para o aperfeiçoamento do regime jurídico de proteção da personalidade, especialmente no que diz respeito à introdução de princípios antropológicos explícitos na regulação da IA, ao fortalecimento dos mecanismos preventivos e à articulação entre os diferentes instrumentos normativos existentes.
